\section{Limiti dei Codici a Blocchi e Codici Perfetti}
Abbiamo già visto che la distanza minima $d_{min}$ di un codice, ne determina le capacità correttive.
È emerso anche il compromesso fondamentale che governa la codifica di canale: per aumentare $d_{min}$ e garantire una maggiore robustezza al rumore, è necessario aumentare il numero di cifre di ridondanza, il che riduce inesorabilmente il tasso informativo ($R = k/n$) e l'efficienza della trasmissione.
Nasce quindi spontanea una domanda cruciale: fissata la lunghezza del blocco $n$ e la distanza minima desiderata $d_{min}$, qual è il \textbf{numero massimo di parole di codice} (e quindi la dimensione $k$ massima) che possiamo "impacchettare" nello spazio vettoriale senza che le loro sfere di decodifica si sovrappongano?

\subsection{Il Limite di Singleton} % Slide 5
Il primo limite alla dimensione dei codici è dato dal Limite di Singleton: è la disuguaglianza più \textit{diretta} e generale che lega i parametri fondamentali di un codice a blocchi.

Sia $C(n,k)$ un codice a blocchi con distanza minima $d_{min}$ su alfabeto di cardinalità $s$, allora:
\[|C|\leq s^{n-d_{min}+1}\]
\textbf{Dimostrazione}: se $|C|$ fosse maggiore del limite, allora \textbf{almeno} due parole di codice $x_1$ e $x_2$ coinciderebbero sulle prime $n-d_{min}+1 $ cifre. Questo comporterebbe che la distanza minima tra le due parole sarebbe
\[d_H(x_1,x_2) \leq n - (n-d_{min}+1) = d_{min}-1 \implies \textbf{ impossibile!}\]
Visualizziamo il ragionamento come analogia del \textit{\textbf{Principio dei cassetti}}:
\begin{enumerate}
    \item \textit{Eliminiamo} le ultime $d_{min} - 1$ cifre da tutte le parole.
    \item rimangono \textit{tronconi} di parole lunghi $n-(d_{min}-1)=n-d_{min}+1$.
    \item le combinazioni diverse (\textit{i cassetti}) per parole lunghe $n-d_{min}+1$ su un alfabeto di cardinalità $s$ sono $s^{n-d_{min}+1}$
    \item ma il codice di partenza ha $s^n$ parole, e $s^k > s^{n-d_{min}+1}$
    
\item questo implica che avremmo due parole di codice in partenza diverse che finiscono nello stesso \textit{cassetto} (ovvero hanno i primi $n-d_{min}+1$ simboli identici). Essendo identiche in questa prima parte, l'unico punto in cui possono differire sono esattamente le $d_{min}-1$ cifre che avevamo \textit{eliminato}. Ma se differiscono al massimo in $d_{min}-1$ posizioni, la loro distanza reciproca sarà minore o uguale a $d_{min}-1$, violando così la definizione stessa di distanza minima del codice!
    
    \item per evitare questo assurdo, si deve imporre che il numero di parole sia minore delle combinazioni possibili
    \item  deve valere $$ |C| \leq s^{n - d_{min} + 1} $$
    
\end{enumerate}


Quindi, limitare la cardinalità del codice o la distanza minima delle sue parole sono due facce della stessa medaglia.


Nel caso specifico dei codici lineari su un alfabeto di cardinalità $s$, il numero delle parole di codice (ovvero il numero totale da codificare) è dato da $|C| = s^k$. Sostituendo questo valore nella disuguaglianza generale del limite di Singleton, si ottiene:
\[ s^k \leq s^{n-d_{min}+1} \]
Applicando il logaritmo in base $s$ (con $s \ge 2$) a entrambi i membri, la disuguaglianza mantiene il proprio verso poiché la funzione logaritmo con base maggiore di 1 è strettamente crescente:
\[ \log_s(s^k) \leq \log_s(s^{n-d_{min}+1}) \]
Dalle proprietà dei logaritmi ($\log_b(b^x) = x$), gli esponenti si "abbassano":
\[ k \leq n - d_{min} + 1 \iff d_{min} \leq n - k + 1\]
Quindi per i codici lineari validi, la distanza minima è dettata dal valore di $k$ scelto.

\subsection{I Limiti di Hamming (Sphere Packing Bound)} % Slide 6-15
Se il Singleton bound è un limite \textit{esistenziale} per i codici lineari, ora analizziamo i limiti legati al \textit{volume geometrico} dello spazio di Hamming, calcolando quante parole di codice possiamo \textit{impacchettare} al suo interno senza che le loro sfere di decodifica si sovrappongano, inficiando così una corretta decodifica.

Prima evidenziamo come sono in relazione il numero di codeword (su alfabeto binario) di un codice e l'\textit{informazione} generata da una sorgente che emette $R$ simboli binari al secondo:
per trasmettere adeguatamente l'informazione generata, se un canale può trasmettere \textbf{un solo simbolo} al secondo, servirà un insieme di $2^{nR}$ parole di codice.

Da ciò si intuisce come per un alto tasso di trasmissione è necessario un numero di parole di codice \textit{grande}.

Ma quante parole di codice possono esserci in un codice affinché si possano correggere configurazioni con al più $t$ errori per parola trasmessa?
In altre parole quanto grande può essere l'insieme di parole del codice, e quindi quanta informazione può passare in un canale, volendo però poter correggere un numero congruo di possibili errori?




\subsubsection{Hamming Upper Bound (Numero di parole di codice)} % Slide 6-11
Consideriamo ogni codeword come il centro della propria \textit{sfera di correzione}, all'interno dello spazio di Hamming. Considerando un codice \textbf{binario}, avremo $M=2^k$ sequenze binarie lunghe $n$, in grado di correggere tutte le configurazioni fino a $t$ errori. Per evitare sovrapposizioni tra sfere/sequenze deve valere la relazione:
\[\textbf{Limite di Hamming:} \space 2^k= 2^{n\frac{k}{n}}=2^{nR} \leq \frac{2^n}{\sum_{i=0}^t\binom{n}{i}}\]
In dettaglio si ha che:
\begin{itemize}
\item $2^n$ è la dimensione dello spazio per \textbf{tutte} le sequenze binarie lunghe $n$.
\item $$\sum_{i=0}^t\binom{n}{i}$$ quantifica quante sequenze differiscono da $w_i$ fino a $t$ posizioni
\end{itemize}

\textbf{Dimostrazione:}
\begin{itemize}
    \item $$w_1,\dots, w_M$$ parole di codice \textbf{binarie}
    \item Ad ogni parola è associata una sfera con \textit{centro} $w_i$, e \textit{raggio} $t$
    \item ognuna delle $M$ \textit{sfere} (partizioni) è l'insieme di \textbf{tutte} le sequenze binarie $v$ tali che $d_H(w_i, v) \leq t$
    \item ogni sequenza che ricade dentro una sfera $i$ dista al massimo $t$ dal suo centro (parola con cui viene decodificata)
    \item ogni sfera \textit{contiene} $1+ \binom{n}{1}+\dots+\binom{n}{t}=\sum_{i=0}^t\binom{n}{i}$ sequenze
    \item  il totale delle sequenze incluse nelle sfere è quindi $M\sum_{i=0}^t\binom{n}{i}$
    \item il numero di sequenze contenute nelle sfere però \textbf{non può eccedere il numero totale di sequenze binarie} in uno spazio di dimensione $n$:
    \[ M\sum_{i=0}^t \binom{n}{i} \leq 2^n
    \implies
    M=2^k\leq \frac{2^n}{\sum_{i=0}^t\binom{n}{i}}\]
\end{itemize}
Il valore di $M$ è un rapporto tra lo spazio e il numero di configurazioni lunghe fino a $t$.

Il limite di Hamming è una \textbf{condizione necessaria ma NON sufficiente} per un codice che corregge fino a $t$ errori.
\begin{esempio}{}
\[t=1,n=4 \implies M\leq \frac{2^n}{(1+n)}= \frac{16}{5}\]
Il numero intero più grande che soddisfa il limite di Hamming è $M=3$. Ma non esistono codici che possano correggere singoli errori con più di 2 parole di lunghezza 4, perché la distanza minima per correggere un errore è $d_{min}=3$
\end{esempio}

\begin{esempio}{}
    Possiamo correggere singoli errori, con parole di lunghezza $5$ con un codice di $4$ parole binarie ($00000,\space 10011,\space 11100,\space 01111$).
\end{esempio}
\begin{esempio}{}
    Se vogliamo correggere errori doppi mantenendo parole di lunghezza $5$, non possiamo usare più di due parole ($00000,\space 11111$).
\end{esempio}


\subsubsection{Hamming Lower Bound (Numero di cifre di controllo)} % Slide 12-15
Con la precedente analisi si è data una limitazione alla dimensione del codice, con questa ulteriore valutazione si impone un limite inferiore al numero $q=n-k$ di cifre di controllo.

\[2^q= 2^{n-k}\geq \sum_{i=0}^t\binom{n}{i}\]
\textbf{Dimostrazione}:
Se un codice a controllo di parità è capace di correggere \textbf{fino a $t$} errori, allora ogni combinazione di \textbf{$t$ o meno} errori deve risultare in un \textbf{singolo correttore} (sindrome distinta).
\begin{itemize}
    \item Il numero di possibili vettori di errore di lunghezza $n$ contenenti esattamente $i$ errori è dato dal coefficiente binomiale $\binom{n}{i}$.
    \item I possibili correttori (sindromi) sono invece $2^{n-k}=2^{q}$.
    \item Per garantire l'unicità di ogni correttore per ogni combinazione di errore deve valere la formulazione del limite di Hamming:
    \[2^{n-k}\geq 1 + \binom{n}{1}+\dots+ \binom{n}{t}=
    \sum_{i=0}^t\binom{n}{i}\]
    \item il termine 1 nella sommatoria è presente nel caso in cui \textbf{non si verifichino} errori (ovvero il vettore di errore nullo, a cui è associata la sindrome nulla);
\end{itemize}
L'Hamming Lower Bound è una \textbf{condizione necessaria MA NON sufficiente} per costruire codici a blocchi che correggano fino a $t$ errori. Potrebbero esistere valori di $q_0$, inteso come il più piccolo intero positivo che soddisfi il Lower Bound, per valori di $n$ e $t$, ma potrebbero \textbf{non} esistere codici validi con $q_0$ cifre di controllo.
\begin{esempio}{Lower Bound soddisfatto, codice impossibile}
\[n=10 \quad t=2 \implies 2^q= 2^{n-k}\geq \sum_{i=0}^t\binom{n}{i}\implies q_0=6\]
Non esistono codici lineari binari che possano correggere un doppio errore con meno di $7$ cifre di controllo, perché la distanza minima per correggere un doppio errore è 5!    
\end{esempio}
Servirà un ulteriore teorema per dare un limite superiore al numero di cifre di controllo.

\paragraph{Hamming Upper e Lower Bound}
Il limite \textit{inferiore} sul \textbf{numero di cifre di controllo} ($q$) e il limite \textit{superiore} sul \textbf{numero delle parole di codice} ($k$) sono, nei fatti, due facce della stessa medaglia.
Partendo dalla formulazione del Lower Bound sulle cifre di controllo $q = n-k$:
\[ 2^{n-k} \geq \sum_{i=0}^t\binom{n}{i} \]
Sfruttando le proprietà delle potenze $2^{n-k} = \frac{2^n}{2^k}$):

\[ \frac{2^n}{2^k} \geq \sum_{i=0}^t\binom{n}{i} \]
A questo punto, moltiplicando entrambi i membri per $2^k$ e dividendo per la sommatoria, ricaviamo esattamente la formula dell'Upper Bound:
\[ 2^k \leq \frac{2^n}{\sum_{i=0}^t\binom{n}{i}} \]
Questo dimostra che limitare lo spazio per le parole di codice o garantire spazio sufficiente per le sindromi d'errore sono lo stesso identico vincolo matematico.


\subsection{Il Limite VGS (Varshamov-Gilbert-Sacks)} % Slide 16-21
Il limite di Hamming definisce una condizione \textit{necessaria} sul numero di cifre di controllo $q$, ma non ci garantisce che un codice con tali parametri sia effettivamente costruibile. 

Adesso introduciamo una condizione \textbf{sufficiente, ma NON necessaria} per la costruibilità del codice.
Un limite \textit{sufficiente} (ma non necessario) per la costruzione di un codice è fornito dal teorema di Varshamov-Gilbert-Sacks (VGS); afferma che:

Un codice binario con parole lunghe $n$ che corregge $t$ errori \textbf{può essere costruito} se il numero delle cifre di controllo soddisfa la seguente relazione:
\[2^q=2^{n-k} > \sum_{i=0}^{2t-1}\binom{n-1}{i}\]
Se per un dato $n$ e un dato $t$, troviamo il più piccolo intero $q_0$ che soddisfa questa relazione, il limite VGS ci dà la certezza matematica che \textbf{esiste} almeno un codice lineare con $q_0$ cifre di controllo in grado di correggere $t$ errori. 

\begin{esempio}{}
    \[n=10,t=2 \implies 2^q > \sum_{i=0}^{2t-1}\binom{n-1}{i} \implies q_0=8\]
    Tuttavia esistono codici che correggono errori doppi con un $q=7$ cifre di controllo.
\end{esempio}

\paragraph{Confronto tra i limiti di Hamming e VGS} % Slide 18 e 21

Applicando il logaritmo in base 2 alle disuguaglianze viste precedentemente, possiamo isolare il numero di cifre di controllo $q$ ed evidenziare il ruolo complementare dei due limiti.

\vspace{0.5cm}
\noindent
\begin{minipage}[t]{0.48\textwidth}
    \centering
    \textbf{Limite di Hamming}\\
    \textit{(Condizione Necessaria)}
    
    \[ q \ge \log_2 \left[ \sum_{i=0}^t\binom{n}{i} \right] \]
    
    Velocità di trasmissione:
    \[ R \le 1 - \frac{1}{n}\log_2 \left[ \sum_{i=0}^t\binom{n}{i} \right] \]
    
    \vspace{0.4cm}
    $\Downarrow$
    \vspace{0.2cm}
    
    \fbox{\begin{minipage}{0.9\textwidth}
        \small \centering
        \textbf{Limite superiore sulla ridondanza minima}: tutti i codici che correggono $t$ errori devono rispettare questo limite.
    \end{minipage}}
\end{minipage}
\hfill
\begin{minipage}[t]{0.48\textwidth}
    \centering
    \textbf{Limite VGS}\\
    \textit{(Condizione Sufficiente)}
    
    \[ q > \log_2 \left[ \sum_{i=0}^{2t-1}\binom{n-1}{i} \right] \]
    
    \vspace{1.75cm} % Spazio per allineare le frecce e i box
    $\Downarrow$
    \vspace{0.2cm}
    
    \fbox{\begin{minipage}{0.9\textwidth}
        \small \centering
        \textbf{Limite inferiore sufficiente sulla ridondanza minima}: esiste almeno un codice che corregge $t$ errori se si rispetta questo limite.
    \end{minipage}}
\end{minipage}
\vspace{0.5cm}

\subsection{Codici Perfetti}
Esiste tuttavia una classe ristretta e molto speciale di codici in cui il limite di Hamming viene soddisfatto con l'\textbf{uguaglianza esatta}:
\[ 2^k \sum_{i=0}^t \binom{n}{i} = 2^n \]
I codici che soddisfano questa equazione sono definiti \textbf{Codici Perfetti}.

Dal punto di vista geometrico, in un codice perfetto le sfere di decodifica di raggio $t$ realizzano una partizione esatta dell'intero spazio di Hamming: riempiono tutto il volume disponibile senza sovrapporsi e \textbf{senza lasciare alcuno "spazio vuoto"}. 

Di conseguenza, \textit{qualsiasi} sequenza binaria $y$ di lunghezza $n$ che il ricevitore possa captare, ricadrà matematicamente dentro una (e una sola) sfera di decodifica. Questo implica che i vettori di errore più pesanti ($wt(e) > t$) non genereranno mai una richiesta di ritrasmissione, ma subiranno inesorabilmente una \textbf{decodifica errata certa}.

\begin{esempio}
I \textbf{Codici di Hamming} (che correggono un singolo errore, $t=1$) sono l'esempio per eccellenza di codici perfetti. Un'altra classe famosa e molto rara in teoria dell'informazione è rappresentata dai \textit{Codici di Golay} (con $n=23,k=12,t=3$). Oltre a questi, esistono solo codici perfetti "banali" (come i codici a ripetizione di lunghezza dispari).
\end{esempio}